Dans l'espace muni d'un repère orthonormé $\left(\text{O}~;~\vec{\imath},~\vec{\jmath},~\vec{k}\right)$ de l'espace, on considère les points suivants

\[A(0~;~0~;~1)~;\quad B(1~;~2~;~3)~;\quad C(3~;~3~;~1)~;\quad E(2~;~-2~;~2)~;\quad F(3~;~0~;~4)\quad \text{et}\quad G(5~;~1~;~2)
\]

\begin{enumerate}
\item \begin{enumerate}
\item Montrer que les points B, C et E ne sont pas alignés.
\item Justifier que le vecteur $\overrightarrow{\mathrm{AF}}$ est normal au plan $(\mathrm{BCE})$.
\item En déduire qu'une équation cartésienne du plan $(\mathrm{BCE})$ est $x+z-4=0$.
\end{enumerate}

\item \begin{enumerate}
\item Montrer que le point G n'appartient pas au plan $(\mathrm{BCE})$.
\item Montrer que les vecteurs $\overrightarrow{\mathrm{BE}}$, $\overrightarrow{\mathrm{BC}}$ et $\overrightarrow{\mathrm{AG}}$ ne sont pas coplanaires.
\item En déduire que la droite $(\mathrm{AG})$ et le plan $(\mathrm{BCE})$ sont sécants.
\end{enumerate}

Pour la suite de l'exercice, on appellera P le point d'intersection de la droite $(\mathrm{AG})$ et du plan $(\mathrm{BCE})$.

\item \begin{enumerate}
\item Montrer qu'une représentation paramétrique de la droite $(\mathrm{AG})$ est :
\[
\left\{
\begin{array}{l}
x=5t\\
y=t\\
z=1+t
\end{array}
\right.\quad t\in \R
\]
\item En déduire les coordonnées du point P.
\item Montrer que le point P est le milieu du segment $[\mathrm{EC}]$.
\end{enumerate}

\item Déterminer l'intersection des plans $(\mathrm{BCE})$ et $(\mathrm{ACG})$.
\end{enumerate}
